\documentclass[a4paper,11pt,final]{article}
% Pour une impression recto verso, utilisez plutôt ce documentclass :
%\documentclass[a4paper,11pt,twoside,final]{article}

\usepackage[english,francais]{babel}
\usepackage[utf8]{inputenc}
\usepackage[T1]{fontenc}
\usepackage[pdftex]{graphicx}
\usepackage{setspace}
\usepackage{hyperref}
\usepackage[french]{varioref}
\usepackage[babel=true]{csquotes}
\usepackage{subfig}
\usepackage{wrapfig}
\usepackage{geometry}
\usepackage{listings}
\usepackage[usenames,dvipsnames,svgnames,table]{xcolor}
\usepackage{float}

\geometry{top=2.7cm,bottom=2.7cm,left=2.7cm,right=2.7cm}
\interfootnotelinepenalty=10000
\widowpenalty=1000
\clubpenalty=1000
\makeatletter
\sloppy
\renewcommand\paragraph{\@startsection{paragraph}{4}{\z@}%
  {-1.5ex\@plus -1ex \@minus -.2ex}%
  {1.5ex \@plus .2ex}%
  {\normalfont\normalsize\bfseries}}
\makeatother

\newcommand{\reporttitle}{Déploiement IPv6 \\ \emph{Adressage, auto-configuration, routage, DHCPv6, mobilité}}     % Titre
\newcommand{\reportauthor}{Émilien \textsc{Kenler}\\ Abdoul \textsc{Khadre Diallo} \\ Clément \textsc{Mercier}\\ Arthur \textsc{Puyou}} % Auteur
\newcommand{\reportsubject}{Déploiement IPv6} % Sujet
\newcommand{\HRule}{\rule{\linewidth}{0.5mm}}
\setlength{\parskip}{1ex} % Espace entre les paragraphes

\lstdefinelanguage{cisco}{
     basicstyle=\small\ttfamily\color{black},
     moredelim=*[l][\color{Gray}]{router},
     moredelim=*[l][\color{Black}]{\#},
     morecomment=[l][\color{Gray}]{!},
     moredelim=**[s][\color{YellowOrange}\ttfamily]{<}{>}
}

\lstnewenvironment{code}[1][]%
{
    \noindent
	\minipage{\linewidth}
     \lstset{
		language=cisco,
         frame=single,
         showstringspaces=false,
         tabsize=5,
         numbers=left,
         firstnumber=1,
         numberstyle=\tiny\color{Blue},
         stepnumber=5,
         extendedchars=false,
         breaklines=true
	 }
}{\endminipage}

\lstnewenvironment{paquet}[1][]%
{
    \noindent
	\minipage{\linewidth}
     \lstset{
         basicstyle=\tiny\ttfamily\color{black},
         frame=single,
         showstringspaces=false,
         tabsize=5,
         numbers=left,
         firstnumber=1,
         numberstyle=\tiny\color{Blue},
         stepnumber=5,
         extendedchars=false,
         breaklines=true
	 }
}{\endminipage}

% Warning avec un doigt
\usepackage{pifont,mdframed}
\newenvironment{warning}
  {\par
    \begin{list}{}{\leftmargin=1cm
                   \labelwidth=\leftmargin}\item[\Large\ding{43}]}
  {\end{list}\par}

% Includes a figure \scalefig{filename}{scale}{legend}
\newcommand{\scalefig}[3]{
  \begin{figure}[H]
    \centering
    \includegraphics[width=#2\columnwidth]{#1}
    \caption{#3}
    \label{#1}
  \end{figure}}
\newcommand{\scalefiguntitled}[2]{
  \begin{figure}[H]
    \centering
    \includegraphics[width=#2\columnwidth]{#1}
    \label{#1}
  \end{figure}}


\hypersetup{
    pdftitle={\reporttitle},%
    pdfauthor={\reportauthor},%
    pdfsubject={\reportsubject},%
    pdfkeywords={rapport} {vos} {mots} {clés}
}

\begin{document}
  \begin{center}


\includegraphics [width=50mm]{logo-utc.pdf} \\[3cm]

\HRule \\[0.2cm]
\begin{spacing}{1.4}
{\huge \bfseries \reporttitle}\\[0.2cm]
\end{spacing}
\HRule \\[0.5cm]
\textsc{\Large SR04 -- A13}\\[2.0cm]

\begin{minipage}[t]{0.3\textwidth}
  \begin{flushleft} \large
    \emph{Étudiants} \\
    \reportauthor
  \end{flushleft}
\end{minipage}
\begin{minipage}[t]{0.6\textwidth}
  \begin{flushright} \large
    \emph{Responsable de l'UV SR04} \\
    M.~Abdelmadjid \textsc{Bouabdallah} \\
  \end{flushright}
\end{minipage}

\vfill

\end{center}
  \cleardoublepage % Dans le cas du recto verso, ajoute une page blanche si besoin
  
  \tableofcontents
  
  \cleardoublepage
  \section{Introduction}

IPv6 est la nouvelle version du protocole IP (couche 3 du modèle OSI), succédant à IPv4, qui est aujourd'hui couramment utilisée.
Elle a été développée pour plusieurs raisons.
\begin{itemize}
  \item À cause de la pénurie d’adresses au niveau mondial, il fallait pouvoir proposer une solution permettant d'augmenter l'espace d'adressage.
  \item Plusieurs champs de l'en-tête IPv4 sont peu ou pas du tout utilisés. IPv6 apporte une simplification du format de l'en-tête.
  \item Les changements de l'en-tête permettent une gestion plus flexible des options et permet l'ajout de nouvelles options dans le futur.
  \item Des fonctionnalités d'authentification et de confidentialité auparavant séparées du protocole IP, sont maintenant directement implémentées dans le protocole.
\end{itemize}
  
Cette étude présente les caractéristiques de cette nouvelle version, puis quelques détails sur son implémentation, notamment sur la partie mobilité.
  

  
  \cleardoublepage  
  \section{Présentation des caractéristiques d’IPv6}
\subsection{Notation des adresses IPv6}
  
\begin{itemize}
  \item La première grande différence est la longueur des adresses qui passe de 32 bits à 128 bits. Le nombre d’adresses augmente donc considérablement.
  \item On passe d’une écriture décimale (\texttt{172.69.31.15}) à une hexadécimale (\texttt{2001:0db8:0000:85a3:0000:0000:ac1f:8001})
  \item Il est permis de ne pas noter les zéros non significatifs. Ainsi l'adresse précédente devient : \texttt{2001:db8:0:85a3:0:0:ac1f:8001}
  \item Des groupes de bits nuls peuvent être omis en gardant les deux-points. Par exemple, dans \texttt{2001:db8:0:85a3::ac1f:8001}, on sait qu'il s’agit de groupes nuls à la place de \texttt{::}. Cette substitution ne peut se faire qu’à un seul endroit, sinon l’adresse devient ambiguë : \texttt{2001:db8::85a3::ac1f:8001} n'est pas valide car on ne sait pas où sont les groupes nuls et comment repartir leur nombre.
  \item Il n’y a plus de masques sous la forme de quatre groupes d'un octet, comme avec IPv4.
Seule la notation CIDR a été conservée.
Elle consiste en un slash (\texttt{/}), suivi de la longueur du préfixe réseau.
Le couple formé par le préfixe réseau et la longueur du préfixe permet de savoir de quel type d’adresse il s’agit.
\end{itemize}

\subsection{Types d’adresses IPv6}

\begin{itemize}
  \item	\texttt{::/128} : adresse non spécifiée, c’est le \emph{this} pour l’attribution d’une adresse.
  \item	\texttt{::1/128} : adresse loopback, la machine elle-même.
  \item	\texttt{fc00::/7} : adresse locale unique, équivalent des adresses privées.
  \item \texttt{fe80::/10} : adresse locale de lien (\textit{link-local address}) pour communiquer avec les appareils à distance \texttt{1}, c’est-à-dire les équipements reliés directement.
On dit que l’adresse est restreinte à un lien.
Nous avons déjà le préfixe : les 64 derniers bits dédiés à l’adresse hôte ont la valeur de l’identifiant d’interface, basée sur l’adresse MAC de la machine.
  \item \texttt{ff00::/8} : adresses multicast.
  \item \texttt{ff02::2} : adresse multicast spécifique permettant de s’adresser à tous les routeurs de distance \texttt{1}.
  \item \texttt{ff02::1} : adresse multicast spécifique permettant de s’adresser à tous les hôtes voisins de distance \texttt{1}.
  \item \texttt{/64} : préfixe d’un sous réseau. La taille d’une sous réseau est fixe en IPv6.
Il y a 48 bits pour le FAI, 16 bits pour le sous réseau et 64 bits pour l’hôte.
\end{itemize}

\subsection{Principe des options}

La plupart des en-têtes des paquets sont désormais optionnelles, elle n'ont plus besoin d'être présentes dans l'en-tête.
Cela permet d’alléger de manière générale la taille d’un paquet tout en permettant une plus grande flexibilité et évolutivité.
Les options de l’en-tête sont stockées sous la forme d’une liste chaînée.

\newpage
\section{Attribution des adresses en IPv6}

\subsection{Auto-configuration}

Il existe différentes méthodes au sein d'IPv6 pour l’attribution d’une adresse IP globale aux équipements.

\begin{itemize}
  \item L'attribution manuelle classique.
  \item L’auto-configuration sans état, en lien avec le NDP.
Il n’y a pas besoin de serveur pour distribuer les adresses, cette dernière est auto-attribuée grâce à NDP.
On a ici un fonctionnement \emph{plug-and-play} au niveau de l’interface réseau.
  \item L’auto-configuration avec état, c’est-à-dire avec un serveur DHCPv6, très semblable à celle sous IPv4.
On garde la main sur les adresses attribuées en donnant un pool d’adresse au serveur DHCP, permettant de centraliser les adresses attribuées.
Nous ne détaillerons pas plus cela. % TODO retirer après avoir fait la partie DHCPv6
\end{itemize}

\subsection{Neighbor Discovery Protocol (NDP)}

Nous allons ici faire le lien entre l’adresse locale de lien (\textit{link-local address}) et l’auto-configuration sans état (\textit{stateless address autoconfiguration}).
Le NDP est une des nouveautés d’IPv6 et c’est ce qui contribue à son attrait.
Ce protocole fournit des services semblables à ARP (Address Resolution Protocol), ICMP (Internet Control and Error Messages) et plus encore.
Il utilise principalement le protocole ICMP qui a été enrichi par rapport à IPv4 : on parle d’ICMPv6.

Voici les différentes fonctions que remplit ce protocole :
\begin{itemize}
  \item La résolution d’adresse, semblable à ARP seulement ici le protocole utilise les messages standards d’ICMPv6.
  \item NUD (Neighbor Unreachability Detection) : c’est une des nouveautés d’IPv6, cela permet d’effacer des tables les équipements proches devenus inaccessibles. Il permet une plus grande adaptabilité pour changer une route au sein d’un routeur si un de ses voisins devient inaccessible.
  \item La configuration automatique d’un équipement grâce à la découverte de ses voisins avec par exemple :
    \begin{itemize}
      \item découverte des routeurs;
      \item découverte des préfixes réseau, qui lui permet de construire son ou ses propres adresses IPv6;
      \item découverte des adresses dupliquées, qui permet d’éviter que l’adresse construite automatiquement n’existe déjà sur le réseau;
      \item découverte des paramètres du lien physique (taille du MTU, nombre de sauts max autorisés...).
    \end{itemize}
  \item	Indication de redirection, envoyée par un routeur s’il connaît une meilleure route en nombre de saut pour atteindre une destination.
\end{itemize}Avec ICMPv6, de nombreuses fonctions sont présentes et de nombreuses sont nouvelles, ce qui enrichit considérablement IPv6.
\newpage
Voici la procédure d’auto-configuration d’une adresse pour un nouvel équipement rattaché à un réseau :
\begin{itemize}
  \item la toute première étape consiste à créer l'adresse locale de lien (\textit{link-local address}) de l’équipement.
  \item une fois l'unicité de cette adresse vérifiée, la machine est en mesure de communiquer avec les autres machines de distance 1. Les machines de distance 1 sont les machines du lien.
  \item pour savoir comment sera attribuée l’adresse unicast globale (l’IP globale), la machine doit chercher à acquérir un message d'annonce (\textit{Router Advertisement}) d’un routeur.
  \item s'il y a un routeur sur le lien, la machine doit appliquer la méthode indiquée par le message d'annonce du routeur :
    \begin{itemize}
      \item l'auto-configuration sans état, la machine utilise sont adresse de lien locale comme adresse unicast globale;
      \item l'auto-configuration avec état (DHCPv6).
    \end{itemize}
  \item en l'absence de routeur sur le lien, la machine doit essayer d'acquérir une adresse unicast globale par la méthode d'auto-configuration sans état. Si la tentative échoue, c'est terminé. Les communications se feront uniquement sur le lien avec l'adresse locale de lien. En effet la machine n'a pas une adresse avec une portée qui l'autorise à communiquer avec des machines autres que celles du lien.
\end{itemize}


  
  \cleardoublepage  
  \section{DHCPv6}

Le DHCPv6 est l’équivalent de la DHCP (IPv4) en IPv6.
C’est un protocole de configuration dynamique.
Il permet d’attribuer les adresses IPv6. 

Il a été mis en place car IPv6 voulait se débarrasser de la dépendance avec DHCP pour configurer les adresses. 
Ainsi, les hôtes peuvent maintenant configurer plusieurs adresses eux-mêmes. 

En plus d’être un protocole de configuration, il joue aussi le rôle de mécanisme d’auto configuration au même titre que NDP. 

DHCPv6 opére dans plusieurs modes :
\begin{itemize}
  \item stateless mode: fourniture de serveurs DNS et d’autres services commes les options pour les téléphones SIP
  \item stateful mode : une machine peut configurer son adresse IP avec DHCPv6
  \item DHCPv6-PD: obtention d’un préfixe réseau par un routeur à partir d’un fournisseur de service
\end{itemize}

\scalefig{images/dhcpv6_1.png}{.7}{Communication client-serveur DHCPv6}

Nous avons les différents types de flags suivants : 
\begin{itemize}
  \item o : indique qu’il y a d’autres informations à trouver dans le serveur 
  \item m: mode géré c’est à dire que le client doit demander une adresse IP et ne pas en configurer une de manière statique (statelessly)
\end{itemize}

Il est aussi à noter que les machines s’identifient au serveur DHCP avec le DUID contrairement à IPv4 où on se connecte avec l’adresse MAC. 
DUID signifie Device Unique IDentifier (identifiant unique de l’appareil). Et chaque participant DHCP en a un. C’est celui-ci qui l’identifie tout au long de la communication.

Enfin, pour communiquer, des messages DHCP sont envoyés; voici  la structure de ces messages: 

\scalefig{images/dhcpv6_2.png}{.5}{Structure des messages DHCPv6}


  
  \cleardoublepage
  \section{Mobilité IPv6 (MIPv6)}


Pour communiquer avec des serveurs, des connexions sont établies et elles sont liées à l’adresse IP.
Le problème lié à la mobilité est que l’adresse IP est liée au réseau sur lequel on se trouve, dès lors que l’on va se déplacer on va changer de réseau et obtenir une nouvelle adresse IP, les connexions actives vont donc être perdues car l’adresse IP avec laquelle elles avaient établi les connexions n’est plus joignable.
Un mécanisme à donc été mis en place au sein d’IPv6 pour résoudre le problème : il faut autoriser un équipement à garder son adresse IP \textbf{même s’il change de réseau}.
Pour cela un équipement est attaché à un réseau, ce sera son \textbf{réseau mère}, qui lui fournira une adresse IP fixe et ce dernier permettra d’assurer la continuité des connexions établies.
On parle de \emph{home address} (HoA) pour l’adresse IP fixe liée au réseau mère (\emph{home network}).
De plus l’équipement aura toujours une adresse IP variable et temporaire liée à sa localisation et au réseau auquel il est directement connecté, on parle de \emph{care-of address} (CoA).
La pierre angulaire de ce système est l’agent mère (HA) situé sur le réseau mère et chargé d’assurer la correspondance entre la HoA et la CoA.

L’IPv6 étant un protocole construit en amont par l’IETF, certains mécanismes pour la mobilité ont directement pu être intégrés dans la conception du protocole.

Voici un petit schéma pour résumer et expliciter tout cela :

\scalefig{images/schema1}{.7}{}

On retrouve le nœud correspondant, qui est le serveur à joindre, le réseau mère qui fait la liaison entre les 2 adresses et le réseau étranger où se trouve physiquement notre équipement.
La HoA se situe donc sur le réseau mère et la CoA se situe sur le réseau étranger.
Voyons maintenant les cas d’utilisations de ce système.

\subsection{Le noeud mobile est dans son réseau mère}

C’est le cas le plus simple, il utilise directement sa \emph{Home Address} et on a un routage classique pour la communication entre le noeud mobile et le serveur.

\scalefig{images/schema2}{.7}{}

\subsection{Le mobile est dans un réseau étranger}

C’est là que cela devient intéressant, le mobile aura donc sa HoA et une adresse étrangère CoA acquise par les mécanismes d’auto-configuration sur le réseau étranger où est physiquement le noeud mobile.
L’agent mère va donc tenir sa table d’association et associer la HoA avec la CoA.
Si un correspondant veut envoyer un paquet au noeud mobile il enverra un paquet avec comme source la sienne et comme destination la \emph{Home Address} du noeud mobile.
Une fois arrivé à l’agent mère le paquet va être routé vers la CoA du noeud mobile.
Pour cela l’agent mère encapsule le paquet qui a donc comme source l’adresse de l’agent mère et comme destination la CoA du mobile.
Une fois arrivé au mobile le paquet est désencapsulé et transmis aux couches supérieurs avec donc comme source et destinations celles du premier paquet.
Au yeux du mobile c’est comme s’il était sur son réseau mère, les connexions ne sont donc plus coupées car même si la CoA change, au yeux du mobile la destination est toujours sa HoA qui ne change pas.
L’encapsulation se fait grâce à l’extension d’en-tête IP-IP d’IPv6, dès lors les paquets sont protégés par IPsec et tunnelé vers la CoA en toute sécurité.

\scalefig{images/schema3}{.7}{}
\scalefig{images/schema4}{.7}{}

\subsection{Mécanisme d’association}

Pour créer ou mettre à jour une association au niveau de l’agent mère, lors d’un changement de réseau du noeud mobile par exemple, un mécanisme spécifique de Mobilité IPv6 est utilisé.
Il s’agit du \emph{Binding Update (BU)}, appelé en français mise à jour d’association.
Ce mécanisme utilise un mécanisme du NDP permettant de savoir qu’un hôte n’est plus accessible.

\scalefig{images/schema5}{.7}{}

Un soucis que pose ce mécanisme est une certaines inefficacité et une lenteur au niveau du routage car il faut obligatoirement passé par le réseau mère, même si le réseau avec lequel on communique est proche de nous.
Un mécanisme d’optimisation du routage a donc été mis en place, mais il n’est pas toujours supporté par le serveur avec lequel on communique.
Il faut donc garder à l’esprit que la mobilité se fait bien sans cela, il ne s’agit que d’une amélioration, pas d’une obligation.

\subsection{Optimisation du routage}

Pour gagner en efficacité dans certains cas, l’optimisation de routage a été mise en place.
Le principe est simple et reprend le mécanisme général avec le réseau mère, seulement ici le noeud correspondant va aussi tenir une table d’association HoA/CoA, comme l'agent mère.
Ainsi par ce biais-là, on garde le même mécanisme, mais le noeud correspondant se substitue à l'agent mère.
Le noeud mobile contacte directement au noeud correspondant, qui va enlever l’en-tête et obtenir un paquet venant de la HoA et à destination du CN, la connexion n’est donc pas brisée.

\scalefig{images/schema6}{.7}{}
\scalefig{images/schema7}{.7}{}

Il faut aussi faire des Binding Update pour garder à jour les tables d’associations.
Le noeud mobile doit donc connaître les agents auxquels il doit envoyer ces mises à jour.
Il a donc une table des agents et lors d’un BU il enverra à son agent mère et au noeud correspondant supportant l’optimisation de routage.
Il doit toujours tenir à jour l’agent mèrei, car certaines connexions ne supportent pas l’optimisation suivant le serveur correspondant, il peut donc y avoir à la fois des connexions optimisées et d’autres non optimisées.

\scalefig{images/schema8}{.7}{}


  
  \cleardoublepage
  \section{Mise en œuvre d'IPv6}

Afin d'illustrer ce projet par une démonstration technique, un routeur \emph{Cisco 1811} nous a été prêté. L'objectif est de créer deux réseaux et de faire basculer un client d'un réseau à un autre, tout en restant joignable sur son adresse IPv6 d'origine.

\subsection{Installation}

Afin de permettre à tous de travailler sur le routeur, le routeur a été installé de manière fixe derrière une \emph{Freebox v6}.
Cette Freebox dispose d'une connexion IPv6 par défaut et chaque abonné dispose de 8 sous-réseaux \texttt{/64}.

Les branchements sont les suivants :
\begin{center}
    \begin{tabular}{|l|l|}
    \hline
    Port                   & Connecté à \\ \hline
    FastEthernet 0 (WAN)   & Freebox    \\ \hline
    FastEthernet 1 (WAN)   & n.c.       \\ \hline
    FastEthernet 2-3 (LAN) & Clients (réseau mère)   \\ \hline
    FastEthernet 4-5 (LAN) & Clients (réseau mobile)   \\ \hline
    FastEthernet 6-9 (LAN) & n.c.   \\ \hline
    \end{tabular}
\end{center}

Pour commencer, il faut se connecter au routeur à l'aide du port console.
La première étape consiste à identifier le routeur en lui donnant un nom :

\begin{code}
router(config)#hostname router
router(config)#ip domain-name emilienkenler.net
\end{code}

Il est alors d'activer le SSH pour pouvoir se connecter à distance :

\begin{code}
router(config)#ip ssh rsa keypair-name sshkey
router(config)#crypto key generate rsa usage-keys label sshkeys \
               modulus 2048
router(config)#ip ssh version 1
router(config)#ip ssh logging events
router(config)#ip ssh time-out 60
router(config)#ip ssh authentication-retries 5
router(config)#service password-encryption
router(config)#enable secret sr04
router(config)#username sr04 password 0 sr04
router(config)#line vty 0 4
router(config-line)#login local
router(config-line)#transport input ssh
\end{code}

\begin{warning}
Le protocole SSH ne semble pas fonctionner en version 2 : le client affiche \texttt{Invalid modulus length}.
La version 1 est donc forcée.
\end{warning}

Pour que les clients puissent joindre Internet, il faut définir la route par défaut en IPv4, qui est la Freebox :

\begin{code}
router(config)#ip route 0.0.0.0 0.0.0.0 192.168.1.254
\end{code}

Il est maintenant possible d'activer l’IPv6 et l’auto-configuration sur l’interface raccordée à la Freebox, dans le but de récupérer une adresse dans le réseau géré par la Freebox :

\begin{code}
router(config)#ipv6 unicast-routing 
router(config)#int FastEthernet 0
router(config-if)#ipv6 enable
router(config-if)#ipv6 address autoconfig default
router#show ipv6 route
S   ::/0 [2/0]
     via FE80::224:D4FF:FE9A:18E0, FastEthernet0
\end{code}

La commande suivante permet de vérifier qu'une adresse a bien été assignée :
\begin{code}
router#sh ipv6 interface FastEthernet 0
  IPv6 is enabled, link-local address is FE80::5A8D:9FF:FE14:8900 
  No Virtual link-local address(es):
  Stateless address autoconfig enabled
  Global unicast address(es):
    2A01:E35:8A38:6870:5A8D:9FF:FE14:8900, subnet is 2A01:E35:8A38:6870::/64 [EUI/CAL/PRE]
      valid lifetime 86004 preferred lifetime 86004
  Joined group address(es):
    FF02::1
    FF02::2
    FF02::1:FF14:8900
  MTU is 1500 bytes
  ICMP error messages limited to one every 100 milliseconds
  ICMP redirects are enabled
  ICMP unreachables are sent
  ND DAD is enabled, number of DAD attempts: 1
  ND reachable time is 30000 milliseconds (using 18502)
  ND advertised reachable time is 0 (unspecified)
  ND advertised retransmit interval is 0 (unspecified)
  ND router advertisements are sent every 200 seconds
  ND router advertisements live for 1800 seconds
  ND advertised default router preference is Medium
  Hosts use stateless autoconfig for addresses.
\end{code}

L'adresse du routeur est donc \texttt{2a01:e35:8a38:6870:5a8d:9ff:fe14:8900}.
Pour éviter de devoir taper cette adresse à chaque connexion, nous avons créé une entrée DNS\footnote{\emph{Domain Name System}} de type \texttt{AAAA} avec le nom \texttt{router.emilienkenler.net} :

\begin{code}
router.emilienkenler.net. 1651  IN      AAAA    2a01:e35:8a38:6870:5a8d:9ff:fe14:8900
\end{code}

Au sein de l'interface de la \emph{Freebox v6}, il est possible de demander le routage de certains réseau vers une adresse qui se trouve derrière la \emph{Freebox}.
Le routeur est alors appelé le \emph{Next Hop} : c'est le point de passage suivant pour tous les paquets adressés à ce sous-réseau.
Ceci va nous permettre d'accéder à chacun des réseaux que nous utilisons pour la démonstration depuis Internet.

Dans notre cas, le réseau par défaut de la \emph{Freebox} est \texttt{2a01:e35:8a38:6870::/64}.
Dans l'interface de la Freebox, on accède donc au menu \emph{Paramètres de la Freebox > Mode avancé > Configuration IPv6}, puis on configure les 4 sous-réseaux qui seront utilisés par le projet pour les faire pointer vers le routeur :

\scalefig{images/freebox}{.7}{Configuration de la Freebox}

Les réseaux seront utilisés de cette manière :

\begin{center}
    \begin{tabular}{|l|l|l|l|}
    \hline
    ID du sous-réseau & Adresse du sous-réseau           & Description                          & Interface \\ \hline
    2                 & \texttt{2a01:e35:8a38:6872::/64} & réseau mère                      & VLAN2     \\ \hline
    3                 & \texttt{2a01:e35:8a38:6873::/64} & Réseau mobile                        & VLAN3     \\ \hline
    4                 & \texttt{2a01:e35:8a38:6874::/64} & \emph{Réservé Wi-Fi}                       & -         \\ \hline
    5                 & \texttt{2a01:e35:8a38:6875::/64} & Réseau extérieur & WAN       \\ \hline
    \end{tabular}
\end{center}

La configuration initiale est maintenant terminée et la console peut être débranchée.

Comme SSH2 ne fonctionne pas, il faut bien penser à forcer la version 1 du protocole pour se connecter au routeur :

\begin{code}
$ ssh -6 -1 sr04@router.emilienkenler.net
\end{code}

\subsection{Auto-configuration IPv6}

Nous allons utiliser le sous-réseau \texttt{2a01:e35:8a38:6875::/64} pour tester la mise en place de l'auto-configuration au niveau du routeur.

Comme il a été délégué au routeur dans la configuration de la Freebox, celle-ci ne l'annonce plus.
Nous pouvons donc configurer ce sous-réseau sur l'interface WAN puis vérifier que les clients connectés au niveau de la Freebox parviennent toujours à le rejoindre.

La configuration se fait ainsi :

\begin{code}
router(config)#int FastEthernet 0
router(config-if)#ipv6 enable
router(config-if)#ipv6 address 2a01:e35:8a38:6875::/64 eui-64
router(config-if)#ipv6 nd prefix 2a01:e35:8a38:6875::/64
\end{code}

La configuration de l'interface montre que celle-ci a déjà récupéré une adresse dans ce sous-réseau :

\begin{code}
router#sh ipv6 interface FastEthernet 0
FastEthernet0 is up, line protocol is up
  IPv6 is enabled, link-local address is FE80::5A8D:9FF:FE14:8900 
  No Virtual link-local address(es):
  Global unicast address(es):
    2A01:E35:8A38:6870:5A8D:9FF:FE14:8900, subnet is 2A01:E35:8A38:6870::/64
    2A01:E35:8A38:6875:5A8D:9FF:FE14:8900, subnet is 2A01:E35:8A38:6875::/64 [EUI]
\end{code}

Un \texttt{ping} depuis l'extérieur montre que cette adresse est bien valide :

\begin{code}
$ ping6 2A01:E35:8A38:6875:5A8D:9FF:FE14:8900
PING 2a01:e35:8a38:6875:5a8d:9ff:fe14:8900 from 2001:41d0:8:9b46::1 eth0: 56 data bytes
64 bytes from 2a01:e35:8a38:6875:5a8d:9ff:fe14:8900: icmp_seq=1 ttl=56 time=32.0 ms
\end{code}

\subsection{Configuration des VLAN}

Afin de séparer les réseaux domicile et mobile entre les différents ports du switch, nous allons utiliser deux VLAN.

Le VLAN2 est utilisé pour le réseau mère.
Il est accessible sur les ports 2 et 3 du switch :

\begin{code}
router(config)#vlan 2
router(config-vlan)#name VLAN2

router(config)#interface FastEthernet 2
router(config-if)#switchport mode access
router(config-if)#switchport access vlan 2

router(config)#interface fastEthernet 3
router(config-if)#switchport mode access
router(config-if)#switchport access vlan 2
\end{code}

Il faut ensuite indiquer l'adresse du sous-réseau correspondant à l'interface et activer l'auto-configuration.

\begin{code}
router(config)#interface vlan 2
router(config-if)#ipv6 enable                              
router(config-if)#ipv6 address 2a01:e35:8a38:6872::/64 eui-64
router(config-if)#ipv6 nd ra interval 20
router(config-if)#ipv6 nd advertisement-interval
router(config-if)#ipv6 nd prefix 2a01:e35:8a38:6872::/64
\end{code}

Une configuration similaire est appliquée pour le VLAN3, qui est utilisé pour le réseau mobile sur les ports 4 et 5 :

\begin{code}
router(config)#vlan 3
router(config-vlan)#name VLAN3

router(config)#interface FastEthernet 4
router(config-if)#switchport mode access
router(config-if)#switchport access vlan 3

router(config)#interface FastEthernet 5
router(config-if)#switchport mode access
router(config-if)#switchport access vlan 3

router(config)#interface vlan 3
router(config-if)#ipv6 enable       
router(config-if)#ipv6 address 2a01:e35:8a38:6873::/64 eui-64
router(config-if)#ipv6 nd ra interval 20
router(config-if)#ipv6 nd advertisement-interval
router(config-if)#ipv6 nd prefix 2a01:e35:8a38:6873::/64
\end{code}

\subsection{Configuration du Wi-Fi}

Nous souhaitions initialement réaliser une démonstration consistant à passer un client d'une connexion filaire à une connexion Wi-Fi.
La conservation de la même IP devrait alors permettre de faire cette bascule sans interruption des connexions.

Cependant, l'interface Wi-Fi du routeur dont nous disposons nécessite l'utilisation d'un bridge qui ne supporte pas l'IPv6.
Il n'est donc pas possible de faire de l'IPv6 en Wi-Fi avec ce routeur.

La configuration du Wi-Fi a donc été abandonnée.
Les commandes pour IPv4 sont toutefois disponibles en annexe à titre de référence.

\subsection{Configuration de Mobile IPv6 sur le routeur}

Afin que les clients puissent utiliser Mobile IPv6 sur les VLAN, le routeur va jouer le rôle de \emph{Home Agent} :

\begin{code}
router(config)#interface vlan 2
router(config-if)#ipv6 mobile home-agent
router(config)#interface vlan 3
router(config-if)#ipv6 mobile home-agent
\end{code}

Quatre commandes sont disponibles pour afficher des détails sur Mobile IPv6 une fois qu'il est activé (elles sont assez verbeuses et les détails ne seront pas reproduits ici) :

\begin{code}
router#sh ipv6 mobile globals
Mobile IPv6 Global Settings:
...

router#sh ipv6 mobile home-agent
Home Agent information for Vlan2
...

router#sh ipv6 mobile binding 
Mobile IPv6 Binding Cache Entries:
...

router#sh ipv6 mobile traffic
MIPv6 statistics:
...
\end{code}

En cas de besoin, on peut activer l'affichage de messages supplémentaires dans les journaux :

\begin{code}
router#debug ipv6 mobile AR-selection
router#debug ipv6 mobile binding-cache
router#debug ipv6 mobile forwarding 
router#debug ipv6 mobile home-agent   
router#debug ipv6 mobile networks    
router#debug ipv6 mobile registrations      
router#debug ipv6 mobile router 
\end{code}

\subsection{Configuration du client}

Il existe plusieurs implémentations de Mobile IPv6 sous Linux.

Nous avons choisi \emph{UMIP} \footnote{Universal Mobile IP for Linux, \url{http://www.umip.org/}}, qui est une implémentation de la RFC6275\footnote{\url{http://tools.ietf.org/html/rfc6275}}.

Une autre alternative est \emph{MIPL} \footnote{Mobile IPv6 for Linux, \url{http://www.mobile-ipv6.org} (plus disponible)}, dont une documentation est encore disponible \footnote{\url{http://tldp.org/HOWTO/Mobile-IPv6-HOWTO/mipv6.html}}, mais dont le site principal, et donc les sources, n'est plus accessible.

\emph{UMIP} nécessite le support du noyau pour pouvoir fonctionner. Sous Fedora, ce support est disponible nativement et le paquet \texttt{mipv6-daemon} peut être installé pour l'installer :

\begin{code}
# yum install mipv6-daemon
\end{code}

On peut alors créer le fichier de configuration \texttt{/etc/mip6d.conf}, qui spécifie que le client joue ici le rôle de \emph{Mobile Node} (\texttt{MN}), ainsi que l'adresse du réseau mère et du \emph{Home Agent} (c'est-à-dire le routeur) :

\begin{code}
NodeConfig MN;
DebugLevel 10;
OptimisticHandoff enabled;
DoRouteOptimizationMN disabled;
MnMaxHaBindingLife 60;

Interface "em1" {
  MnIfPreference 1;
}

MnHomeLink "em1" {
  HomeAgentAddress 2A01:E35:8A38:6872:5A8D:9FF:FE14:8900;
  HomeAddress 2A01:E35:8A38:6872::1/64;
}
\end{code}

Il faut ensuite redémarrer le démon pour prendre en compte cette configuration :

\begin{code}
# systemctl restart mip6d.service
\end{code}

\subsection{Analyse du traffic}

Maintenant que tout est configuré, on peut réaliser la démonstration.
L'objectif est de connecter un client dans le réseau mère, puis de le connecter dans le réseau mobile et de vérifier qu'il est toujours joignable sur son adresse d'origine.
Pour vérifier cela, on utilisera un tiers qui envoie des \texttt{ping} vers l'adresse dans le réseau mère.
Le tiers pourra se trouver dans un réseau quelconque.

Les paquets affichés dans ce document sont issus de l'analyse avec \emph{Wireshark} de la capture \texttt{capture1.pcapng}.
La première colonne de chaque ligne correspond au numéro de séquence dans cette capture.

\begin{warning}
Une première remarque importante à faire concerne la gestion de l'état des interfaces VLAN par le routeur.

Conformément à une note sur le site de Cisco\footnote{\url{https://supportforums.cisco.com/docs/DOC-2032}}, une interface VLAN n'est disponible que si au moins un périphérique est connecté à une interface dans ce VLAN.

Ainsi, si le client mobile de test est le seul périphérique connecté au switch et qu'on le branche sur un port du VLAN mobile, les ports 2 et 3 ne sont pas connectés.
Le switch va alors désactiver le VLAN2 et passer l'interface à l'état \texttt{up/down} au lieu de \texttt{up/up}.

Si on regarde l'état du Mobile IPv6 sur le routeur dans ces conditions, on constate effectivement qu'aucune adresse n'est détectée sur l'interface VLAN2 :
\begin{code}
router#sh ipv6 mobile home-agent
Home Agent information for Vlan2
  Configured:
    FE80::5A8D:9FF:FE14:8900
    preference 0 lifetime 1800
      No global addresses within prefix

  No Discovered Home Agents
Home Agent information for Vlan3
  Configured:
    FE80::5A8D:9FF:FE14:8900
    preference 0 lifetime 1800
      global address 2A01:E35:8A38:6873:5A8D:9FF:FE14:8900/64

  No Discovered Home Agents
\end{code}

Une alternative consiste à utiliser le VLAN dans lequel se trouve la Freebox, mais celle-ci va alors répondre aux trames liées à Mobile IPv6, ce qui perturbe le client.

Il faut donc veiller à laisser un périphérique connecté en permanence dans le réseau mère.
\end{warning}

On déroule les étapes suivantes :

\begin{enumerate}
\item Activation de l'interface sur le client A
\item Ping de B vers A, puis arrêt du ping
\item A est débranché du VLAN3 et branché sur le VLAN2
\item Ping de B vers A puis arrêt du ping
\item Démarrage de \texttt{mip6d} sur le client A
\item Ping de B vers A en continu
\item A est débranché du VLAN2 et branché sur le VLAN3
\item Le ping doit toujours fonctionner
\end{enumerate}

\subsubsection{Auto-configuration}

% J'ai pas trop compris ce qu'il se passe... Je dirais que effd reçoit un Advertisement dès le début, donc il fait un Neighbor Solicitation et c'est réglé direct ?

Après activation de l'interface réseau sur le client A, celui-ci s'auto-configure en émettant un \texttt{Router Solicitation} :

\begin{paquet}
7	1.386	fe80::a00:27ff:fe9e:effd	ff02::2	ICMPv6	70	Router Solicitation from 08:00:27:9e:ef:fd
\end{paquet}

Le routeur répond avec un \texttt{Router Advertisement} :

\begin{paquet}
10	3.166	fe80::5a8d:9ff:fe14:8900	ff02::1	ICMPv6	126	Router Advertisement from 58:8d:09:14:89:00
\end{paquet}

\subsubsection{Test dans chaque réseau}

Le premier test consiste à envoyer un ping de B vers A dans le réseau mobile.

\begin{paquet}
28	74.503	2a01:e35:8a38:6872:226:b0ff:feed:e94a	2a01:e35:8a38:6873:a00:27ff:fe9e:effd	ICMPv6	70	Echo (ping) request id=0x0450, seq=2, hop limit=63 (reply in 29)
29	74.503	2a01:e35:8a38:6873:a00:27ff:fe9e:effd	2a01:e35:8a38:6872:226:b0ff:feed:e94a	ICMPv6	70	Echo (ping) reply id=0x0450, seq=2, hop limit=64 (request in 28)
\end{paquet}

L'adresse de A dans le réseau mobile est \texttt{2a01:e35:8a38:6873:a00:27ff:fe9e:effd}, ce qui est bien une adresse dans le VLAN3.
L'adresse de B est \texttt{2a01:e35:8a38:6872:226:b0ff:feed:e94a}.
Elle est dans le VLAN2 et ne changera pas au cours de la démonstration.

Nous débranchons ensuite A du VLAN3 pour le placer dans le VLAN2, attendons que son IP soit fixée et relançons un ping vers son adresse dans le réseau mère.

\begin{paquet}
83	121.933	2a01:e35:8a38:6872:226:b0ff:feed:e94a	2a01:e35:8a38:6872:a00:27ff:fe9e:effd	ICMPv6	70	Echo (ping) request id=0x0452, seq=0, hop limit=64 (reply in 84)
84	121.933	2a01:e35:8a38:6872:a00:27ff:fe9e:effd	2a01:e35:8a38:6872:226:b0ff:feed:e94a	ICMPv6	70	Echo (ping) reply id=0x0452, seq=0, hop limit=64 (request in 83)
\end{paquet}

Cette étape préliminaire permet de vérifier que la communication fonctionne bien entre les deux réseaux avant de mettre en œuvre l'IPv6.

\subsubsection{Mobilité}

En laissant A dans le VLAN2, c'est-à-dire son réseau mère, le démon \texttt{mip6d} est lancé :

\begin{code}
# systemctl start mip6d.service
\end{code}

Ce service envoie alors un \texttt{Binding Update} de A vers le routeur pour indiquer qu'il est dans son réseau mère.
Le routeur répond avec un \texttt{Binding Acknowledgment}.

\begin{paquet}
364	290.533	2a01:e35:8a38:6872::1	2a01:e35:8a38:6872:5a8d:9ff:fe14:8900	MIPv6	110	Binding Update
365	290.536	2a01:e35:8a38:6872:5a8d:9ff:fe14:8900	2a01:e35:8a38:6872::1	MIPv6	94	Binding Acknowledgement
\end{paquet}

Le ping de B vers A est alors démarré.
Le ping est maintenant effectué vers l'adresse qui a été définie dans la configuration du \texttt{mip6d}, soit \texttt{	2a01:e35:8a38:6872::1}.

\begin{paquet}
437	339.593	2a01:e35:8a38:6872:226:b0ff:feed:e94a	2a01:e35:8a38:6872::1	ICMPv6	70	Echo (ping) request id=0x0457, seq=0, hop limit=64 (reply in 438)
438	339.593	2a01:e35:8a38:6872::1	2a01:e35:8a38:6872:226:b0ff:feed:e94a	ICMPv6	70	Echo (ping) reply id=0x0457, seq=0, hop limit=64 (request in 437)
\end{paquet}

Ensuite, A est débranché du VLAN2 et branché sur un port du VLAN3.

\begin{warning}
Durant la phase de découverte, les ping de B vers A ne sont pas visibles.
Ceci provient du fait que la capture est réalisée depuis A.
\end{warning}

Après quelques secondes, A avertit le routeur de son réseau mère avec un \texttt{Binding Update} :

\begin{paquet}
485	355.962	2a01:e35:8a38:6872::1	2a01:e35:8a38:6872:5a8d:9ff:fe14:8900	MIPv6	110	Binding Update
...
    Alternate care-of address: 2a01:e35:8a38:6873:a00:27ff:fe9e:effd
\end{paquet}

Ce paquet contient l'adresse \emph{Care Of} dans le réseau mobile : \texttt{2a01:e35:8a38:6873:a00:27ff:fe9e:effd}

La réponse revient à nouveau sous forme de \texttt{Binding Acknowledgement} :

\begin{paquet}
488	356.965	2a01:e35:8a38:6872:5a8d:9ff:fe14:8900	2a01:e35:8a38:6872::1	MIPv6	94	Binding Acknowledgement
\end{paquet}

Immédiatement, les ping retransmis par le réseau mère sont visibles, et la réponse de A est transmise :

\begin{paquet}
493	357.601	2a01:e35:8a38:6872:226:b0ff:feed:e94a	2a01:e35:8a38:6872::1	ICMPv6	110	Echo (ping) request id=0x0457, seq=18, hop limit=63 (reply in 494)
494	357.601	2a01:e35:8a38:6872::1	2a01:e35:8a38:6872:226:b0ff:feed:e94a	ICMPv6	110	Echo (ping) reply id=0x0457, seq=18, hop limit=64 (request in 493)
\end{paquet}

On notera que la valeur \texttt{hop limit} passe de $64$ à $63$, car le paquet effectue un saut supplémentaire par le routeur du réseau mère avant d'arriver à A.

\subsection{Sécurisation}

À partir de la partie précédente, il doit être évident que l'utilisation de Mobile IPv6 doit être sécurisée.
En effet, il paraît très simple de réaliser un faux \texttt{Binding Update} pour faire rediriger tout le traffic d'une IP tierce vers une IP que l'on contrôle.

La sécurisation de Mobile IPv6 passe par l'utilisation d'IPSec, qui permet d'assurer l'authenticité et l'intégrité des paquets.
La RFC4877\footnote{\url{http://tools.ietf.org/html/rfc4877}} décrit un tel usage.

Cisco a également développé des méthodes d'authentification, et notamment un système de clé pré-partagées (SPI), qui est décrit dans la RFC4285\footnote{\url{http://tools.ietf.org/html/rfc4285}}.
Cette méthode est simple à mettre en œuvre, mais elle n'est malheureusement pas implémentée dans UMIP.



  \cleardoublepage
  \section{Conclusion}

Cette étude nous a permis d'acquérir de bonnes connaissances de base sur IPv6, qui va devenir très rapidement omniprésent.
Mobile IPv6 est une fonctionnalité supplémentaire apportée par IPv6.
Elle montre que le protocole permet une construction modulaire : la sécurité de Mobile IPv6 peut ainsi être implémentée de différentes manières, tout en conservant un protocole de base compatible.

Cependant, cette modularité est aussi une faiblesse : Cisco et UMIP n'ont pas choisi d'implémenter les mêmes techniques pour la sécurisation, ce qui pose des soucis d'interopérabilité.
Cette modularité fait penser à l'authentification sur les réseaux Wi-Fi, que nous avons étudiée en SR06.
La RFC5778\footnote{\url{http://tools.ietf.org/html/rfc5778}} décrit d'ailleurs une méthode pour s'authentifier avec de l'EAP en Mobile IPv6, un protocole qui est couramment utilisé sur les réseaux Wi-Fi aujourd'hui.

Enfin, 10 ans après la publication de la première RFC sur Mobile IPv6, son support reste relativement limité : une des implémentations sous Linux a été abandonnée et l'utilisation sous les dernières variantes de Debian nécessite une compilation de noyau.
De plus, le processus nécessite une configuration manuelle qui réserve pour l'instant son utilisation à des cas très précis.

Mobile IPv6 nécessite donc encore quelques développements avant d'être utilisable à grande échelle, mais apporte une fonctionnalité précieuse à IPv6.
Le nécessaires sera donc très probablement réalisé avec le déploiement d'IPv6 dans les années qui viennent.


%                                                               %
%                        \                                      %
%                         \                                     %
%                          \\                                   %
%                           \\                                  %
%                            >\/7                               %
%                        _.-(6'  \                              %
%                       (=___._/` \                             %
%                            )  \ |                             %
%                           /   / |                             %
%                          /    > /                             %
%                         j    < _\                             %
%                     _.-' :      ``.                           %
%                     \ r=._\        `.                         %
%                    <`\\_  \         .`-.                      %
%                     \ r-7  `-. ._  ' .  `\                    %
%                      \`,      `-.`7  7)   )                   %
%                       \/         \|  \'  / `-._               %
%                                  ||    .'                     %
%                                   \\  (                       %
%                                    >\  >                      %
%                                ,.-' >.'                       %
%                               <.'_.''                         %
%                                 <'                            %
%                                                               %

  \cleardoublepage
  \section{Annexe}

\subsection{Configuration du WiFi}

Bien que nous n'ayons pu faire fonctionner, l'IPv6 en WiFi, voici la configuration utilisée pour obtenir une connexion WiFi opérationnelle en IPv4.

Nous avons choisi de créer un vlan différent des autres pour le WiFi.
Nous utiliserons le vlan 4 et le réseau WiFi sera visible (\emph{guest-mode}).

\begin{code}
router(config)#vlan 4
router(config-vlan)#name VLAN04

router(config)#dot11 ssid wifi
router(config-ssid)#guest-mode
router(config-ssid)#vlan 4
router(config-ssid)#authentication open
\end{code}

Une fois le SSID créé, nous pouvons activer l'interface radio.

\begin{code}
router(config)#interface dot11Radio 0
router(config-if)#ssid WIFI
router(config-if)#station-role root
router(config-if)#no shutdown
router(config-if)#antenna receive diversity
router(config-if)#antenna transmit diversity
router(config-if)#no dot11 extension aironet
\end{code}

Il faut ensuite créer le bridge et configurer l'interface qui correspond à la bande de fréquences WiFi que nous voulons utiliser.

\begin{code}
router(config)#bridge irb
router(config)#interface vlan 4
router(config-if)#bridge-group 1
router(config)#interface bvi 1
router(config)#ip address 10.0.1.1 255.255.255.0
router(config)#interface dot11Radio 0.1
router(config-subif)#description WIFI VLAN04
router(config-subif)#encapsulation dot1Q 4 native
router(config-subif)#no cdp enable
router(config-subif)#bridge-group 1
\end{code}

Pour terminer la configuration et rendre le WiFi fonctionnel, il faut configurer le DHCP.

\begin{code}
router(config)#ip dhcp excluded-address 10.0.1.1 10.0.1.50
router(config)#ip dhcp pool dpool1
router(dhcp-config)#network 10.0.1.0 255.255.255.0
router(dhcp-config)#default-router 10.0.1.1
router(dhcp-config)#import all
\end{code}

Il est maintenant possible de se connecter sur le réseau portant le nom \enquote{wifi}.
Celui-ci est ouvert et non protégé.



  \cleardoublepage
  \listoffigures
\end{document}
